\chapter{Lezione 23}
\label{chap:lezione_23} % Etichetta per riferirsi al capitolo

\begin{flushright}
\textit{Data: 27/05/2025}
\end{flushright}
\textit{Bibliografia: 
Most Probable Path of the Ornstein-Uhlenbeck process. Fluctuations and Onsager relations [Puglisi-Sarracino-Vulpiani, ch 3]. Shannon and Gibbs Entropy. Introduction to Stochastic Thermodynamics.}



\section{Traiettorie Ottimali}

Vogliamo calcolare il percorso ottimale che collega due punti.
Poiché stiamo lavorando con un'equazione differenziale stocastica, la soluzione generica $\varphi(t)_{\eta}$ è una traiettoria casuale che dipende dalla realizzazione del rumore $\eta(t)$.
Il percorso ottimale è quello che massimizza la probabilità di transizione $P(\varphi_{1},t_{1}|\varphi_{0},t_{0})$, che rappresenta la probabilità di trovare il sistema nello stato $\varphi_{1}$ al tempo $t_{1}$, dato che si trovava nello stato $\varphi_{0}$ al tempo $t_{0}$.

Questa probabilità può essere espressa attraverso un integrale di cammino (path integral) sulle possibili traiettorie $\varphi(t)$ e sui campi ausiliari $\hat{\varphi}(t)$:
\begin{equation}
P[\varphi_{1},t_{1}|\varphi_{0},t_{0}]=\int_{\varphi_{0}=\varphi(t_{0})}^{\varphi_{1}=\varphi(t_{1})}\mathcal{D}\varphi\mathcal{D}\hat{\varphi}e^{-\frac{S[\varphi,\hat{\varphi}]}{D}}
\end{equation}
L'azione dinamica corrispondente è:
\begin{equation}
S[\varphi,\hat{\varphi}]=\int_{t_{0}}^{t_{1}}ds\mathcal{L}[\varphi,\hat{\varphi}]
\end{equation}
Il Lagrangiano è dato da:
\begin{equation}
\mathcal{L}[\varphi,\hat{\varphi}]=-\hat{\varphi}^{2}+\hat{\varphi}[\partial_{t}\varphi+V^{\prime}(\varphi)]
\end{equation}
Qui, $V'(\varphi)$ rappresenta la derivata del potenziale $V(\varphi)$ rispetto a $\varphi$.

\noindent Consideriamo un potenziale a doppia buca.
Se scegliamo come stato iniziale $\varphi_0$ la posizione di un minimo del potenziale e come stato finale $\varphi_1$ la posizione dell'altro minimo, il sistema deve superare una barriera di potenziale. Il passaggio da un minimo all'altro può avvenire tramite due tipi di processi:
\begin{itemize}
    \item \textbf{Uphill:} Il sistema si muove contro la forza derivante dal potenziale (serve "rumore" o energia esterna).
    \item \textbf{Downhill:} Il sistema si muove seguendo la forza derivante dal potenziale (processo "gratuito" o spontaneo).
\end{itemize}

Estendendo l'intervallo temporale da $t_0 \rightarrow -\infty$ e $t_1 \rightarrow +\infty$, l'eccitazione che porta il sistema da un minimo all'altro è chiamata istantone. L'istantone è una soluzione classica delle equazioni del moto, ovvero una traiettoria che minimizza l'azione funzionale.

Nel limite $D \rightarrow 0$ (dove $D$ è un parametro legato all'intensità del rumore o alla temperatura), la probabilità di transizione può essere approssimata tramite il metodo del punto di sella:
\begin{equation}
P(\varphi_{1},t_{1}|\varphi_{0},t_{0})\simeq \exp \Biggl\{-\frac{S[\varphi_{sp},\hat{\varphi}_{sp}]}{D} \Biggr\}
\end{equation}
dove $(\varphi_{sp},\hat{\varphi}_{sp})$ sono le traiettorie del punto di sella che soddisfano le equazioni del moto ottenute variando l'azione $S$:
\begin{equation}
\frac{\delta S}{\delta\varphi}\Bigg|_{SP}=0 \quad  \quad , \quad  \quad\frac{\delta S}{\delta\hat{\varphi}}\Bigg|_{SP}=0
\end{equation}
Da queste condizioni si ottengono le seguenti equazioni del moto di Ornstein-Uhlenbeck:

\begin{tcolorbox}[colback=colorA!5, colframe=colorA, arc=3mm, boxrule=1.2pt]
\begin{equation}
\begin{cases} 
\partial_{t} \hat{\varphi}_{SP}  &= \hat{\varphi}_{SP} V^{\prime\prime}(\varphi_{SP})     \quad \quad(i)\\  
\partial_{t} \varphi_{SP}  &= - V^{\prime}(\varphi_{SP}) +2 \hat{\varphi}_{SP}   \quad \quad(ii)   
\end{cases}
\end{equation}
\end{tcolorbox}

\subsection{Caso Semplice: Oscillatore Armonico}
Consideriamo il potenziale di un oscillatore armonico:
\begin{equation}
V(\varphi)=\frac{k\varphi^{2}}{2}
\end{equation}
Le derivate del potenziale sono:
\begin{equation}
V^{\prime}(\varphi)=k\varphi \implies -V^{\prime}(\varphi)=-k\varphi
\end{equation}
\begin{equation}
V^{\prime\prime}(\varphi)=k
\end{equation}
Sostituiamo nelle equazioni del moto $(i)$ e $(ii)$ :
\begin{equation}
\begin{cases} 
\partial_{t} \hat{\varphi}_{SP}  &= k\hat{\varphi}_{SP}  \quad \quad(i) \\  
\partial_{t} \varphi_{SP}  &= -k\varphi_{SP} +2 \hat{\varphi}_{SP}   \quad \quad(ii)   
\end{cases}
\end{equation}
Risolviamo queste equazioni.
Dalla prima equazione per $\hat{\varphi}(t)$:
\begin{equation}
\hat{\varphi}_{SP}(t)=\hat{\varphi}_{0}e^{kt}
\end{equation}
Sostituiendo questa soluzione nella seconda equazione per $\varphi(t)$:
\begin{equation*}
(ii) \quad \quad  \partial_{t} \varphi_{SP}  = -k\varphi_{SP} +2 \hat{\varphi}_{0}e^{kt}
\end{equation*}
Questa è un'equazione differenziale lineare del primo ordine. La soluzione generale si ottiene sommando la soluzione dell'omogenea associata e una soluzione particolare dell'equazione completa.
Risolvendo (ponendo $t_0=0$ senza perdita di generalità):
\begin{equation}
\varphi(t)=\varphi_{0}e^{-kt}+2\int_{0}^{t}ds\hat{\varphi}(s)e^{-k(t-s)}
\end{equation}
Sostituiendo $\hat{\varphi}(s)=\hat{\varphi}_{0}e^{ks}$:
\begin{equation}
\varphi(t)=\varphi_{0}e^{-kt}+2\hat{\varphi}_{0}\int_{0}^{t}dse^{ks}e^{-k(t-s)} = \varphi_{0}e^{-kt}+2\hat{\varphi}_{0}e^{-kt}\int_{0}^{t}dse^{2ks}
\end{equation}
Calcolando l'integrale: $\int_{0}^{t}dse^{2ks} = \left[\frac{e^{2ks}}{2k}\right]_0^t = \frac{e^{2kt}-1}{2k}$
Quindi:
\begin{equation}
\varphi(t)=\varphi_{0}e^{-kt}+2\hat{\varphi}_{0}e^{-kt}\frac{e^{2kt}-1}{2k} = \varphi_{0}e^{-kt}+\frac{\hat{\varphi}_{0}}{k}(e^{kt}-e^{-kt})
\end{equation}
Utilizzando la definizione di seno iperbolico $\text{senh}(x) = \frac{e^x - e^{-x}}{2}$:
\begin{equation}
\varphi(t)=\varphi_{0}e^{-kt}+\frac{2\hat{\varphi}_{0}}{k}\text{senh}(kt) \label{eq:sol_generale_armonico}
\end{equation}
Questa è la soluzione generale.
Imponiamo le condizioni al contorno: $\varphi(t_1) = \varphi_1$
\begin{equation}
\varphi_{1}=\varphi(t_{1})=\varphi_{0}e^{-kt_{1}}+\frac{2\hat{\varphi}_{0}}{k}\text{senh}(kt_{1})
\end{equation}
Da questa equazione possiamo ricavare $\hat{\varphi}_0$:
\begin{equation}
\hat{\varphi}_{0}=\frac{k}{2}\frac{\varphi_{1}-\varphi_{0}e^{-kt_{1}}}{\text{senh}(kt_{1})}
\end{equation}
Sostituiendo $\hat{\varphi}_{0}$ nell'espressione (\ref{eq:sol_generale_armonico}):
\begin{equation}
\varphi(t)=\varphi_{0}e^{-kt}+\frac{\varphi_{1}-\varphi_{0}e^{-kt_{1}}}{\text{senh}(kt_{1})}\text{senh}(kt)
\end{equation}
Questa traiettoria è esatta per l'oscillatore armonico nel contesto di questo formalismo.

\section{Termodinamica Stocastica}
La connessione tra probabilità e entropia è data dalla formula di Boltzmann $S=k_{B}\ln W$, dove $W$ è il numero di microstati. Questo suggerisce una forma per la probabilità di una fluttuazione:
\begin{equation}
P(\text{fluttuazione}) \sim e^{\frac{S}{k_{B}}}
\end{equation}
Consideriamo:
\begin{itemize}
    \item Un set di $M$ osservabili/variabili macroscopiche, che arrangiamo in un vettore: $\vec{\alpha} = (\alpha_{1},\dots,\alpha_{M})$.  Un esempio potrebbe essere la magnetizzazione in un sistema di Ising: $m=\frac{1}{N}\sum_{i}s_{i}$ (somma di variabili aleatorie).
    \item Queste variabili macroscopiche $\underline{\alpha}$ sono una funzione delle configurazioni microscopiche $\underline{x}$ del sistema: $\underline{\alpha} = \underline{g}(\underline{x})$. Per un sistema di particelle, $\underline{x}$ può rappresentare le posizioni e i momenti $(q,p)$ di tutte le particelle.
    \item Parametri di controllo: pressione, volume, temperatura, numero di particelle, ecc. Questi parametri definiscono l'ensemble statistico.
\end{itemize}
Quindi, abbiamo $\alpha_m = g_m(\underline{x})$ per $m=1,\dots,M$.
Per un sistema di particelle in contatto con un bagno termico (ensemble canonico), la probabilità di una configurazione microscopica $\underline{x}$ (che può contenere le posizioni $\underline{q}$ e i momenti $\underline{p}$) è data da:
\begin{equation}
P(\underline{x},\underline{p}) = P(\underline{x}, N, V, \beta) = \frac{e^{-\beta H[\underline{x}]}}{Z_{\beta}^{N,V}}
\end{equation}
dove $H[\underline{x}] $ è l'Hamiltoniana del sistema, $\beta = (k_B T)^{-1}$ e $Z_{\beta}^{N,V}$ è la funzione di partizione canonica:
\begin{equation}
Z_{\beta}^{N,V}=\int\frac{d\underline{p}d\underline{q}}{\hbar^{3N}N!}e^{-\beta H[\underline{p},\underline{q}]}
\end{equation}
Siamo interessati alla distribuzione di probabilità $P(\underline{\alpha})$ del nostro set di osservabili.

\noindent Procediamo in maniera formale:

Se $\underline{x}=\{S_{i}\}_{i=1}^{N}$ rappresenta le variabili microscopiche ($(q,p)$ in questo caso), la probabilità delle variabili macroscopiche $\underline{\alpha}$ è ottenuta marginalizzando la distribuzione di probabilità microscopica:
\begin{equation}
P(\underline{\alpha})=\int d\underline{x}\prod_{m=1}^{M}\delta[\alpha_{m}-g_{m}(\underline{x})]\rho(\underline{x},\underline{p})
\end{equation}
dove $\rho(\underline{x},\underline{p})$ è la densità di probabilità canonica che dipende dai parametri di controllo $\underline{P}$

\noindent Le funzioni delta assicurano che consideriamo solo le configurazioni microscopiche $\underline{x}$ che corrispondono ai valori specifici delle osservabili $\alpha_m$.
Per un sistema di particelle:
\begin{equation} 
\rho(\underline{x},\underline{P})=\rho(\underline{x},N,V,\beta) = \frac{e^{-\beta H[\underline{x}]}}{Z(\underline{P})} 
\end{equation}
L'energia libera di Helmholtz $F(\underline{P})$ è definita como:
\begin{equation}
F(\underline{P})=-\frac{1}{\beta}\log Z(\underline{P}) \implies Z(\underline{P})=e^{-\beta F(\underline{P})}
\end{equation}
Quindi la densità di probabilità microscopica può essere scritta come:
\begin{equation}
\rho(\underline{x},\underline{P})=\frac{e^{-\beta H[\underline{x}]}}{Z(\underline{P})}=e^{\beta F(\underline{P})-\beta H[\underline{x}]}
\end{equation}
Sostituiendo in $P(\underline{\alpha})$:
\begin{equation}
P(\underline{\alpha})=e^{\beta F(\underline{P})} \int d\underline{x} e^{-\beta H[\underline{x}]} \prod_{m=1}^{M}\delta[\alpha_{m}-g_{m}(\underline{x})]
\end{equation}
L'integrale può essere interpretato come una funzione di partizione $Z(\underline{\alpha},\underline{P})$ per un valore fissato di $\underline{\alpha}$, $Z(\underline{\alpha},\underline{P}) = e^{-\beta F(\underline{\alpha}|\underline{P})}$, dove $F(\underline{\alpha}|\underline{P})$ è un'energia libera condizionata per le variabili $\underline{\alpha}$.
Quindi, la probabilità $P(\underline{\alpha})$ può essere scritta come:
\begin{equation}
P(\underline{\alpha}) = e^{-\beta[F(\underline{\alpha}|\underline{P})-F(\underline{P})]}
\end{equation}
Identifichiamo la variazione di entropia $\delta S(\underline{\alpha})$ come:
\begin{equation}
\frac{\delta S(\underline{\alpha})}{k_B} = -\beta[F(\underline{\alpha}|\underline{P})-F(\underline{Pctx})]
\end{equation}
Se consideriamo $k_B =1$ si ha:
\begin{equation}
\delta S(\underline{\alpha})= -\frac{1}{T}[F(\underline{\alpha}|\underline{P})-F(\underline{P})]
\end{equation}
Da cui segue:
\begin{equation}
P(\underline{\alpha})=C \cdot \exp\left\{\frac{\delta S(\underline{\alpha})}{k_{B}}\right\}
\end{equation}
Siamo interessati a fluttuazioni attorno all'equilibrio. I valori medi (o più probabili) delle osservabili sono i valori di equilibrio $\langle\underline{\alpha}\rangle=\underline{\alpha}^{*}$.

Sviluppiamo $\delta S(\underline{\alpha})$ in serie di Taylor attorno al valore di equilibrio $\underline{\alpha}^{*}$ fino al secondo ordine (il termine lineare è nullo perché $\underline{\alpha}^{*}$ è un massimo per l'entropia):
\begin{equation}
\delta S(\underline{\alpha})=-\frac{1}{2}\delta\underline{\alpha}^{T}G\delta\underline{\alpha}+O(\delta\alpha^{3})
\end{equation}
dove $\delta\underline{\alpha}=\underline{\alpha}-\langle\underline{\alpha}\rangle=\underline{\alpha}-\underline{\alpha}^{*}$.

La matrice $G$ è definita come la derivata seconda (negativa) dell'entropia, valutata all'equilibrio:
\begin{equation}
G_{ij}=-\frac{\partial^{2}S}{\partial\alpha_{i}\partial\alpha_{j}}\Bigg|_{\underline{\alpha}^{*}}
\end{equation}
Questa forma quadratica per l'entropia porta a una distribuzione di probabilità Gaussiana per le fluttuazioni $\delta\underline{\alpha}$:
\begin{align}
P(\underline{\alpha}) &= C' \exp\left\{-\frac{1}{2k_{B}}\delta\underline{\alpha}^{T}G\delta\underline{\alpha}\right\} \nonumber  \\
&= \sqrt{\frac{(2\pi k_B)^M}{\det G}}\exp\left\{-\frac{1}{2k_{B}}\delta\underline{\alpha}^{T}G\delta\underline{\alpha}\right\}
\end{align}
Poiché i parametri $\underline{\alpha}$ fluttuano gaussianamente attorno all'equilibrio, vogliamo scrivere un'equazione di Langevin per queste variabili.

\noindent Definiamo le affinità $X_m$ coniugate alle variabili $\delta\alpha_m$:
\begin{equation}
X_{m} \equiv \frac{\partial(\delta S)}{\partial(\delta\alpha_{m})}
\end{equation}
Se $\delta S(\underline{\alpha})=-\frac{1}{2}\delta\underline{\alpha}^{T}G\delta\underline{\alpha}$, allora:
\begin{equation}
X_{m} = -\sum_j G_{mj}\delta\alpha_{j}
\end{equation}
Dalla teoria delle fluttuazioni all'equilibrio, si hanno le seguenti relazioni (teorema di equipartizione generalizzato o relazioni di Onsager):
\begin{equation}
\langle\delta\alpha_{j}X_{m}\rangle=-k_{B}\delta_{jm}
\end{equation}
E la covarianza delle fluttuazioni è legata all'inversa della matrice $G$:
\begin{equation}
\langle\delta\alpha_{i}\delta\alpha_{j}\rangle=k_{B}G_{ij}^{-1}
\end{equation}
Per semplicità, trasliamo l'origine in modo che $\underline{\alpha}^{*}=0$, quindi $\delta\underline{\alpha}=\underline{\alpha}-\underline{\alpha}^{*}=\underline{\alpha}$

\noindent Possiamo postulare un'equazione di Langevin lineare multivariata (processo di Ornstein-Uhlenbeck) per le fluttuazioni $\underline{\alpha}$:
\begin{equation}
d\underline{\alpha}=-M\underline{\alpha}dt+Bd\underline{W}
\end{equation}
dove $M$ è la matrice di rilassamento, $B$ è una matrice che descrive l'ampiezza del rumore, e $d\underline{W}$ è un vettore di incrementi di processi di Wiener indipendenti (rumore bianco Gaussiano), con $\langle dW_i(t) \rangle = 0$ e $\langle dW_i(t) dW_j(t') \rangle = \delta_{ij} \delta(t-t') dt$.

\noindent Le matrici $M$ e $B$ sono costanti.

\noindent Da un processo di Ornstein-Uhlenbeck, possiamo derivare la corrispondente equazione di Fokker-Planck per la densità di probabilità $P(\underline{\alpha},t)$:
\begin{equation}
\partial_{t}P(\underline{\alpha},t)= \sum_{j,m} M_{jm}\partial_{\alpha_j}[\alpha_{m}P(\underline{\alpha},t)] + \frac{1}{2}\sum_{j,m} Q_{jm}\frac{\partial^{2}P}{\partial\alpha_{j}\partial\alpha_{m}}
\end{equation}
dove $Q_{jm}$ è la matrice di diffusione, data da $Q = BB^T$
Avendo ridefinito $\delta\underline{\alpha}=\underline{\alpha}$ , si ha:
\begin{equation}
\sigma = \underline{\alpha} \cdot\underline{\alpha}=k_{B}G^{-1}
\end{equation}
Con $\sigma_{ij} = \alpha_i \alpha_j$

\noindent Se $\underline{\alpha} = \begin{pmatrix} x \\ y \end{pmatrix}$, allora $\underline{\alpha} \cdot\underline{\alpha} = \begin{pmatrix} x^2 & xy \\ yx & y^2 \end{pmatrix}$

\noindent Ora definiamo l'equazione di Lyapunov. L'equazione di Lyapunov stazionaria lega la matrice di covarianza $\sigma$ alla matrice di drift $M$ e alla matrice di diffusione $Q=BB^T$ per un processo di Ornstein-Uhlenbeck stazionario:
\begin{equation}
\tilde{Q} = BB^T = M\sigma + \sigma M^T
\end{equation}
Consideriamo l'evoluzione della matrice di covarianza:

$\underline{\alpha} \equiv \underline{\alpha}(t)$

$\underline{\alpha}' \equiv \underline{\alpha}(t+dt) = \underline{\alpha}-M\underline{\alpha}dt+Bd\underline{W}$

\noindent Calcoliamo $\langle\alpha_{i}^{\prime}\alpha_{j}^{\prime}\rangle$:

$\alpha_i' = \alpha_i - \sum_l M_{il}\alpha_l dt + \sum_l B_{il}dW_l$

$\alpha_j' = \alpha_j - \sum_m M_{jm}\alpha_m dt + \sum_m B_{jm}dW_m$
\begin{equation*}
\langle\alpha_{i}^{\prime}\alpha_{j}^{\prime}\rangle = \langle (\alpha_i - \sum_l M_{il}\alpha_l dt + \sum_l B_{il}dW_l) (\alpha_j - \sum_m M_{jm}\alpha_m dt + \sum_m B_{jm}dW_m) \rangle
\end{equation*}
Mantenendo i termini fino a $dt$:

$\langle \alpha_i \alpha_j \rangle - \sum_l M_{il} \langle \alpha_l \alpha_j \rangle dt - \sum_m M_{jm} \langle \alpha_i \alpha_m \rangle dt + \sum_{l,m} B_{il} B_{jm} \langle dW_l dW_m \rangle + O(dt^{3/2})$

Usando $\langle dW_l dW_m \rangle = \delta_{lm} dt$ e $\langle \alpha_k dW_l \rangle = 0$:

\begin{equation}
\sigma_{ij}(t+dt) = \sigma_{ij}(t) - (M\sigma(t))_{ij}dt - (\sigma(t)M^T)_{ij}dt + (BB^T)_{ij}dt
\end{equation}

Quindi l'equazione differenziale per $\sigma(t)$ è $\frac{d\sigma}{dt} = -M\sigma - \sigma M^T + Q$

\noindent All'equilibrio, $\frac{d\sigma}{dt}=0$, il che porta all'equazione di Lyapunov: $M\sigma + \sigma M^T = Q$

\noindent Introduciamo la matrice di Onsager $L$:
\begin{equation}
L = MG^{-1}
\end{equation}
Poichè $\underline{\alpha} = G^{-1} \underline{X}$ segue $X_{m} \equiv \frac{\partial(S)}{\partial(\alpha_m)}$

\noindent Sostituiamo $\underline{\alpha}$ nella relazione $d\underline{\alpha}=-M\underline{\alpha}dt+Bd\underline{W}$ :
\begin{equation}
d\underline{\alpha} = -L\underline{X}dt + Bd\underline{W} 
\end{equation}
Imponiamo che la dinamica sia "time reversal" perché siamo all'equilibrio, e lì vale il bilancio dettagliato.

\noindent Una conseguenza del bilancio dettagliato per osservabili generiche $F, G$ è:
$\langle F[x(t)]G[x(0)]\rangle = \langle G[x(0)]F[x(t)]\rangle$

\noindent Per le variabili $\alpha$ , le correlazioni temporali soddisfano:
\begin{equation}
\langle \underline{\alpha}(t)\underline{\alpha}(t')\rangle = \langle \underline{\alpha}(t')\underline{\alpha}(t)\rangle 
\end{equation}
Consideriamo $t'=t+dt$:
\begin{equation} 
\alpha_{j}(t+dt)=\alpha_{j}(t)-\sum_m L_{jm}X_{m}dt+\sum_m B_{jm}dW_{m} 
\end{equation}
\begin{align} 
\langle\alpha_{i}(t)\alpha_{j}(t+dt)\rangle &= \langle\alpha_{i}(t) [ \alpha_{j}(t)-\sum_m L_{jm}X_{m}dt+\sum_m B_{jm}dW_{m} ] \rangle \nonumber\\ 
&= \langle\alpha_{i}(t)\alpha_{j}(t)\rangle - \sum_m L_{jm}\langle\alpha_{i}(t)X_{m}(t)\rangle dt  
\end{align}
Assumendo $\langle \alpha_i(t) dW_m(t) \rangle = 0$
\begin{align} 
\langle\alpha_{i}(t+dt)\alpha_{j}(t)\rangle &= \langle [ \alpha_{i}(t)-\sum_m L_{im}X_{m}dt+\sum_m B_{im}dW_{m} ] \alpha_{j}(t) \rangle \nonumber \\ 
&= \langle\alpha_{i}(t)\alpha_{j}(t)\rangle - \sum_m L_{im}\langle X_{m}(t)\alpha_{j}(t)\rangle dt
\end{align}
Affinché $\langle\alpha_{i}(t)\alpha_{j}(t+dt)\rangle = \langle\alpha_{j}(t)\alpha_{i}(t+dt)\rangle$, e usando $\langle \alpha_i X_m \rangle = -k_B \delta_{im}$:

$-L_{jm}\langle\alpha_{i}X_{m}\rangle = -L_{im}\langle X_{m}\alpha_{j}\rangle$
$-L_{ji}(-k_B) = -L_{ij}(-k_B)$

\noindent Questo implica $L_{ji} = L_{ij}$ ovvero che la matrice di Onsager è simmetrica.
\begin{equation}
L=MG^{-1} \implies G^{-1}=\frac{1}{k_{B}}\sigma
\end{equation}
e quindi:
\begin{equation}
M\sigma = (M \sigma)^T 
\end{equation}

\section{Entropia e Teoria dell'Informazione}
L'entropia è un modo per quantificare l'incertezza. Data una distribuzione di probabilità $P_i$ per variabili discrete, o $P(x)$ per variabili continue.
L'entropia di Shannon è definita come:
\begin{equation}
H[P]=-\sum_{i}p_{i}\log p_{i}
\end{equation}
Per il caso continuo: $H[P]=-\int P(x)\log P(x) dx$

\noindent Questa definizione è consistente con l'entropia termodinamica. Per un sistema all'equilibrio:
\begin{equation}
S=k_{B}H[P_{eq}]=-k_{B}\sum_{\{c\}}P_{eq}[c]\log P_{eq}[c]
\end{equation}
dove $\{c\}$ rappresenta le configurazioni microscopiche del sistema.

\noindent L'energia libera è data da $F=-\frac{1}{\beta}\log Z \implies Z=e^{-\beta F}$

\noindent La probabilità di equilibrio di una configurazione è:
\begin{equation}
P_{eq}[c]=\frac{e^{-\beta E[c]}}{Z}=e^{\beta (F-E[c])}
\end{equation}
Sostituiendo nell'espressione dell'entropia di Shannon:
\begin{align}
H[P_{eq}]&=-\sum_{\{c\}}P_{eq}[c] \beta (F-E[c]) \nonumber \\ 
&= -\beta F \sum P_{eq}[c] + \beta \sum P_{eq}[c]E[c] \nonumber \\ 
& =-\beta F + \beta \langle E \rangle = \beta (\langle E \rangle - F)
\end{align}
Poiché in termodinamica $S = (\langle E \rangle - F)/T$ si ha:
\begin{equation}
H[P_{eq}]=\frac{S}{k_{B}}
\end{equation}
Vogliamo ora capire cosa succede poco fuori dall'equilibrio (Termodinamica Stocastica)
Consideriamo un sistema descritto da un'equazione di Langevin del tipo:
\begin{equation}
\dot{x}=\mu F(x)+\eta
\end{equation}
Un esempio è un colloide in un bagno termico.
Abbiamo:
\begin{itemize}
    \item La traiettoria stocastica $x_{\eta}(t)$
    \item La densità di probabilità $P(x,t)$
\end{itemize}
Se la forza deriva da un potenziale, $F(x)=-V^{\prime}(x)$, allora per $t \rightarrow \infty$, la distribuzione di probabilità $P(x,t)$ tende alla distribuzione di equilibrio $P_{\infty}(x)=P_{eq}(x)$

\noindent Introduciamo una dipendenza da un parametro di controllo esterno $\lambda$, che può variare nel tempo: $F(x) \rightarrow F(x,\lambda)$. Ad esempio, $\lambda$ potrebbe essere un campo magnetico esterno.

\noindent La forza $F(x,\lambda)$ può essere decomposta in una parte conservativa (derivante da un potenziale) e una parte non conservativa:
\begin{equation}
F(x,\lambda)=-\partial_{x}V(x,\lambda)+f(x,\lambda)
\end{equation}
dove $f(x,\lambda)$ è il termine di forza non conservativa.

\noindent Quando $f \neq 0$ o $\lambda$ varia nel tempo, il sistema può raggiungere uno stato stazionario di non equilibrio:
\begin{equation*}
P(x,\lambda,t) \xrightarrow{t\rightarrow\infty} P_{\infty}(x,\lambda) \neq P_{eq}(x,\lambda)
\end{equation*}
Chiaramente, se $f=0$ e $\lambda=\text{costante}$, e le condizioni al contorno imposte sono compatibili con correnti nulle, allora il sistema raggiunge l'equilibrio:
\begin{equation*}
P(x,t,\lambda) \xrightarrow{t\rightarrow\infty} P_{eq}(x,\lambda)
\end{equation*}
dove:
\begin{equation}
P_{eq}(x,\lambda)=\exp\left\{-\frac{1}{T}[V(x,\lambda)-F(\lambda)]\right\}
\end{equation}
e l'energia libera è:
\begin{equation}
F(\lambda)=-T\log\int dx e^{-V(x,\lambda)/T}
\end{equation}
In un contesto di termodinamica stocastica:
\begin{itemize}
    \item Attraverso la variazione del parametro di controllo $\lambda(t)$ e/o la presenza della forza non conservativa $f$, possiamo fare lavoro $W$ sul sistema.
    \item Questo lavoro $W$ verrà (parzialmente o totalmente) dissipato come calore $Q$.
    \item Ci sarà una produzione di entropia $S_{\text{prod}}$.
\end{itemize}